\section{Related work}
\label{sec:related}
% Every sentence that describes another paper has a row in paper/RELATED_WORK_AUDIT.md.

\paragraph{Edge-level detection on AMLworld.}
AMLworld~\citep{altman2023} is a family of synthetic transaction datasets with laundering
labels on transactions and a temporal 60/20/20 split; detectors on it are compared by
minority-class F1 over transactions. Multi-GNN adapts message passing to directed
multigraphs~\citep{egressy2024}; the Graph Feature Preprocessor feeds streaming subgraph-pattern
features to gradient-boosted trees~\citep{blanusa2024}; FraudGT is a graph
transformer~\citep{lin2024fraudgt}; MEGA-GNN adds two-stage multi-edge
aggregation~\citep{bilgi2024}; SALT-GNN adds degree-aware statistical aggregation before
attention and reports $0.716 \pm 0.018$ HI-Small F1 for its transformer variant (code not
released at the time of writing)~\citep{losavio2026salt}; ExSTraQt reports the highest HI-Small
F1 we found, with tree models on graph features computed from expanding time
windows~\citep{tariq2026exstraqt}. Reported HI-Small F1 for these detectors runs from 0.57
(Multi-GIN) to 0.81 (ExSTraQt on FraudGT's split ratios), under slightly different split ratios;
FraudGT itself reports 0.697 for the base model and 0.764 with port numbering and ego
IDs~\citep{lin2024fraudgt}. All are evaluated per transaction.
% source: paper/RELATED_WORK_AUDIT.md rows 5, 7, 8, 8a, 8b (Egressy T2; FraudGT T2/T3; SALT-GNN T2; ExSTraQt T2 and code) We use Multi-GINe as the reference detector, GFP features with
XGBoost as a weaker control and MEGA-GNN as a stronger one.

\paragraph{Group- and sub-network-level detection.}
Table~\ref{tab:related-comparison} lists the group-level work we found and how each defines
its evaluation. \citet{chai2023} learn to propose the laundering sub-network behind a
transaction and evaluate on private AliPay data. \citet{rico2026} learn account classification
and sub-network discovery jointly; on HI-Small their unit is the sub-network around a target
account, scored by node-wise precision and recall against that account's true sub-network, so
no matching between predicted and true groups is needed. \citet{tian2025} mine laundering
groups collaboratively across institutions and report account-level metrics on AMLworld; we
score their method under our protocol (Section~\ref{sec:results-expansion}).
\citet{tariq2026redirect} propose a turnover-weighted recall over the injected patterns of
the large IBM releases. \citet{wang2026clue} recover a laundering group from an analyst's clue
on HI-Small and HI-Medium, taking connected components of all laundering accounts as the
groups. GARG-AML scores accounts for smurfing on HI-Small, with account labels derived from
transaction labels, and uses FlowScope~\citep{li2020flowscope} as a
baseline~\citep{deprez2025garg}. On blockchain data, Elliptic2 provides labelled subgraphs for
subgraph classification under a random split~\citep{bellei2024,song2024revtrack}.
The closest prior work to our decoupling result is \citet{malik2026queues}, who show that
review queues built from transaction-level and from actor-level scores overlap only partly
under the same review budget (mean Jaccard 0.374 at a 1\,\% budget under temporal evaluation). That study uses the Elliptic++ Bitcoin dataset and random-forest classifiers,
and its unit is the individual actor address, not a group; it does not consider rings or
AMLworld.

% Group-level money-laundering evaluation in prior work. Every cell is verified against the
% paper (or, where marked "code", its released code); "n.v." = not verified, listed in
% paper/RELATED_WORK_AUDIT.md. "not stated" = read, and the paper does not say.
\begin{table*}[t]
\centering
\small
\caption{How group-level money-laundering work defines its evaluation. ``Scope rule'': whether
the paper states how a group that crosses a split or evaluation-window boundary is scored.
``Tail'': whether the high-base-rate last days of the AMLworld data are treated specially.}
\label{tab:related-comparison}
\begin{tabular}{@{}p{2.3cm}p{1.7cm}p{2.0cm}p{3.1cm}p{2.8cm}p{2.1cm}p{1.2cm}p{1.2cm}@{}}
\toprule
Work & Data & Unit & Ground-truth group & Metric / matching & Split & Scope rule & Tail \\
\midrule
\citet{chai2023} & private AliPay & sub-network around a transaction & n.v. & n.v. & n.v. & n.v. & n/a \\
\citet{tian2025} & AMLworld HI, LI; AMLSim; private & account flags from mined sets & not stated for AMLworld & node-level P/R/F1; no group matching & none (unsupervised) & not stated & not stated \\
\citet{rico2026} & HI-Small (and two others) & ego sub-network of an account & paper: connected, laundering accounts only; code: 2-hop ball of laundering accounts within one day & node-wise P/R per target account; no matching & stated 70/15/15 temporal; code: days 1--6 / 7--8 / 9--10 & not stated & code: days $>$10 dropped \\
\citet{tariq2026redirect} & IBM large releases; Libra & per-account community & injected patterns & turnover-weighted recall; best-matched flow, no threshold & none (unsupervised) & not stated & not stated \\
\citet{deprez2025garg} & HI-Small, LI-Large; synthetic & account & none (account labels from edge share) & AUC-ROC, AUC-PR, P@K & random 70/30 & not stated & not stated \\
\citet{wang2026clue} & HI-Small, HI-Medium & members recovered from a clue & connected components of all laundering accounts & recall of members; AUC/AP & not stated & not stated & not stated \\
\citet{bellei2024} & Elliptic2 & given subgraph (classification) & labelled subgraphs & F1, PR-AUC, ROC-AUC & random 80/10/10 & n/a & n/a \\
\citet{song2024revtrack} & Elliptic2 & subgraph; sender--receiver link & labelled subgraphs & F1, PR-AUC; HR, NDCG & random 80/10/10 & n/a & n/a \\
\midrule
This paper & HI-Small, HI-Medium & laundering attempt & one attempt of the generator's pattern file & one-to-one greedy, Jaccard $\ge$ 0.5 (0.3, 0.7 swept) & day split, primary window & stated, four conventions & quantified \\
\bottomrule
\end{tabular}
\end{table*}


Group-level evaluation on AMLworld therefore exists. The works do not share a unit (ego
sub-network, mined set, clue-anchored group, injected pattern, laundering attempt), a ground
truth, a matching rule or a split, and none of the papers states how a group that crosses
the evaluation window is scored. Our contribution is not the task but a protocol that makes
those choices explicit, and a measurement of how much each is worth.

\paragraph{Benchmark artefacts.}
The high laundering rate of AMLworld's last days has been noted by practitioners working with
the data,\footnote{For example \url{https://github.com/Kluey/AML-Transaction-Detection}, which
restricts evaluation to the earlier days for this reason.} and the released code of one
group-level method keeps only the first ten days of HI-Small~\citep{rico2026}. We are not
aware of a published quantification of the effect, of its reproduction across releases, or
of its consequences for ring-level metrics and seed budgets.

\paragraph{Statistics.}
We use bootstrap intervals, paired permutation tests and Holm's sequentially rejective
correction~\citep{holm1979}; the explanation layer described in the appendix uses
TreeSHAP~\citep{lundberg2017,lundberg2020} and Fleiss' $\kappa$~\citep{fleiss1971,landis1977}.
% source: explainability/report/xai_chapter.md §2.1
